\section*{Bab \ref{ch:g8:ice-on-earth} --- Es di Bumi}

\begin{solution}{exo:g8:ice-on-earth:1}
(a) Daerah akumulasi. (b) Daerah ablasi. (c) Moncong \omterm{def:g8:ice-on-earth:glacier}{gletser}. (d) \omterm{def:g8:ice-on-earth:glacial-landforms}{Morena}
(terminal).
\end{solution}

\begin{solution}{exo:g8:ice-on-earth:2}
(a) Salah: \omterm{def:g8:ice-on-earth:glacier}{gletser} terbuat dari \omterm{def:g4:clouds-rain-snow:precipitation}{salju} yang jatuh di daratan lalu berubah
menjadi es; air \omterm{def:g2:water-on-earth:ocean}{laut} yang membeku adalah \omterm{def:g8:ice-on-earth:sea-ice}{es laut}. (b) Benar. (c) Salah:
esnya tetap mengalir menuruni lereng; moncongnya mundur karena es mencair
lebih cepat daripada mengalirnya. (d) Benar.
\end{solution}

\begin{solution}{exo:g8:ice-on-earth:3}
(a) \omterm{def:g8:ice-on-earth:glacial-landforms}{Lembah berbentuk U}. (b) Batu eratik. (c) \omterm{def:g8:ice-on-earth:ice-sheet}{Paparan es}. (d) \omterm{def:g8:ice-on-earth:permafrost}{Permafrost}.
\end{solution}

\begin{solution}{exo:g8:ice-on-earth:4}
(a) Antarktika dan Greenland. (b) Antarktika, sekitar \qty{58}{m}
(Greenland: sekitar \qty{7.4}{m}). (c) Sekitar sepersepuluh (10~\%).
\end{solution}

\begin{solution}{exo:g8:ice-on-earth:5}
(a) Sekitar \num{7.05}~juta~\unit{km^2} pada tahun 1979 dan
\num{3.57}~juta~\unit{km^2} pada tahun 2012. (b) $7.05 - 3.57 = 3.48$:
sekitar \num{3.5}~juta~\unit{km^2}. (c) $3.48 \div 7.05 \approx 0.49$:
sekitar separuhnya. (d) Satu tahun saja bergantung pada \omterm{def:g1:watching-the-weather:weather}{cuaca} musim panas
tahun itu: tahun 2012 berada jauh di bawah garis. Garis tren memakai semua
tahun, sehingga satu musim panas yang tidak biasa hanya sedikit sekali
mengubahnya.
\end{solution}

\begin{solution}{exo:g8:ice-on-earth:6}
(a) $2004 - 1941 = 63$ tahun. (b) $12000 \div 63 \approx 190.5$: sekitar
\qty{190}{m} setahun. (c) $2021 - 1984 = 37$ tahun; $3200 \div 37 \approx
86$: sekitar \qty{86}{m} setahun. (d) \omterm{def:g8:ice-on-earth:glacier}{Gletser} Muir, lebih dari dua kali
lebih cepat.
\end{solution}

\begin{solution}{exo:g8:ice-on-earth:7}
Di \emph{daerah akumulasi}, tinggi di \omterm{def:g1:landscapes-around-us:landform}{gunung}, \omterm{def:g4:clouds-rain-snow:precipitation}{salju} yang turun setiap
tahun lebih banyak daripada yang mencair. \omterm{def:g4:clouds-rain-snow:precipitation}{Salju} baru mengubur dan menekan
\omterm{def:g4:clouds-rain-snow:precipitation}{salju} lama, yang berubah menjadi \emph{firn} setelah sekitar satu tahun,
lalu, lebih dalam lagi dan jauh kemudian, menjadi es yang padat. Karena
beratnya sendiri, es itu \emph{mengalir} perlahan menuruni \omterm{def:g1:landscapes-around-us:landform}{lembah}. Lebih
ke bawah, di \emph{daerah ablasi}, es mencair lebih cepat daripada
turunnya \omterm{def:g4:clouds-rain-snow:precipitation}{salju}; es yang sampai di moncong adalah es yang terbentuk jauh di
atas, sering kali sudah lama sekali.
\end{solution}

\begin{solution}{exo:g8:ice-on-earth:8}
(a) Di gelas B: es yang mencair mengalir ke dalam air dan menaikkannya; di
gelas A ketinggiannya tetap seperti semula. (b) Gelas A adalah \omterm{def:g8:ice-on-earth:ice-sheet}{gunung es}
(es yang sudah mengapung di \omterm{def:g2:water-on-earth:ocean}{laut}); gelas B adalah \omterm{def:g8:ice-on-earth:ice-sheet}{lapisan es} (es di
daratan). (c) \omterm{def:g8:ice-on-earth:ice-sheet}{Paparan es} menahan \omterm{def:g8:ice-on-earth:glacier}{gletser-gletser} di belakangnya. Ketika
\omterm{def:g8:ice-on-earth:ice-sheet}{paparan es} itu pecah, \omterm{def:g8:ice-on-earth:glacier}{gletser-gletser} ini bergerak lebih cepat dan
menuangkan lebih banyak es daratan ke \omterm{def:g2:water-on-earth:ocean}{laut}, dan itulah yang menaikkan
\omterm{def:g4:world-physical-map:altitude}{permukaan laut}.
\end{solution}

\begin{solution}{exo:g8:ice-on-earth:9}
(a) \omterm{def:g8:ice-on-earth:glacial-landforms}{Lembah berbentuk U}, yang diukir oleh \omterm{def:g8:ice-on-earth:glacier}{gletser}. (b) \omterm{def:g8:ice-on-earth:glacier}{Gletser} membawanya
dari jauh di hulu \omterm{def:g1:landscapes-around-us:landform}{lembah}, tempat terdapat granit, lalu meninggalkannya
ketika esnya mencair: itulah batu eratik. (c) \omterm{def:g8:ice-on-earth:glacial-landforms}{Morena} lateral: tingginya
menunjukkan seberapa tebal \omterm{def:g8:ice-on-earth:glacier}{gletser} itu dahulu; es pernah mencapai
setidaknya setinggi itu di sisi \omterm{def:g1:landscapes-around-us:landform}{lembah}.
\end{solution}

\begin{solution}{exo:g8:ice-on-earth:10}
(a) $7.4 \div 10 = \qty{0.74}{m}$. (b) $58 \div 100 = \qty{0.58}{m}$.
(c) $58 \div 7.4 \approx 7.8$: sekitar 8 kali lebih banyak.
\end{solution}

\begin{solution}{exo:g8:ice-on-earth:11}
(a) $7.21 \times 10^6 = \qty{7210000}{km^2}$ dan $4.65 \times 10^6 =
\qty{4650000}{km^2}$. (b) $7.21 - 4.65 = \num{2.56}\text{ juta }\unit{km^2}$,
yaitu \qty{2560000}{km^2}. (c) $2.56 \div 3 \approx 0.85$: sekitar
\num{0.85}~juta~\unit{km^2} per dasawarsa, sedikit lebih besar daripada
garis tren, yang \num{0.76}~juta~\unit{km^2}; rata-rata setiap dasawarsa
masih membawa naik turunnya sepuluh musim panasnya sendiri, tetapi kedua
cara membaca grafik itu menunjukkan kehilangan besar yang sama. (d) Luas
sebaran tidak dapat turun di bawah nol, dan garis tren hanya menggambarkan
tahun-tahun yang diukur: apa yang terjadi selanjutnya bergantung pada
bagaimana \omterm{def:g5:weather-and-climate:climate}{iklim} Arktik berubah, sehingga garis itu bukan sebuah
prakiraan.
\end{solution}

\begin{solution}{exo:g8:ice-on-earth:12}
(a) $120 \div 1 = 120$: garis pantai berada sekitar \qty{120}{km} lebih
jauh ke arah \omterm{def:g2:water-on-earth:ocean}{laut}. (b) Air itu telah menguap dari \omterm{def:g2:water-on-earth:ocean}{samudra}, jatuh sebagai
\omterm{def:g4:clouds-rain-snow:precipitation}{salju} di \omterm{def:g3:continents-and-oceans:continent}{benua-benua}, lalu tetap di sana sebagai es dari lapisan-lapisan es
yang besar. (c) \omterm{def:g8:ice-on-earth:glacial-landforms}{Morena} terminal terbentuk di ujung \omterm{def:g8:ice-on-earth:glacier}{gletser}. \omterm{def:g8:ice-on-earth:glacial-landforms}{Morena}
terminal di sebuah \omterm{def:g1:landscapes-around-us:landform}{dataran} yang jauh dari \omterm{def:g4:world-physical-map:range}{pegunungan} mana pun menunjukkan
bahwa dahulu ada \omterm{def:g8:ice-on-earth:ice-sheet}{lapisan es} yang mencapai \omterm{def:g1:landscapes-around-us:landform}{dataran} itu, pada masa ketika es
menutupi jauh lebih banyak daratan daripada sekarang.
\end{solution}

\begin{solution}{pb:g8:ice-on-earth:1}
\textbf{1.} $1080 - 0 = \qty{1080}{m}$.

\textbf{2.} 1960--1970: \qty{180}{m}; 1970--1980: \qty{150}{m};
1980--1990: \qty{190}{m}; 1990--2000: \qty{180}{m}; 2000--2010:
\qty{180}{m}; 2010--2020: \qty{200}{m}. Paling sedikit pada 1970--1980,
paling banyak pada 2010--2020.

\textbf{3.} Tidak: \omterm{def:g8:ice-on-earth:glacier}{gletser} yang menyusut tetap mengalir menuruni lereng.
Moncongnya mundur karena es mencair lebih cepat daripada datangnya.

\textbf{4.} \omterm{def:g8:ice-on-earth:glacial-landforms}{Morena} terminal: es membawa puing \omterm{def:g3:stones-pebbles-sand:rock}{batuan} menuruni \omterm{def:g1:landscapes-around-us:landform}{lembah} dan
menjatuhkannya di moncongnya, tempat es mencair, selama moncong itu
berada di sana.

\textbf{5.} $1080 \div 6 = \qty{180}{m}$ per dasawarsa.

\textbf{6.} $180 \div 10 = \qty{18}{m}$ setahun.

\textbf{7.} $190 \div 18 \approx 10.6$: sekitar 11 kali lebih cepat.

\textbf{8.} Lebih sedikit: \omterm{def:g8:ice-on-earth:glacier}{gletser} itu kehilangan lebih banyak es di
daerah ablasinya daripada yang diperolehnya di daerah akumulasinya,
sehingga moncongnya mundur.

\textbf{9.} Berbentuk U: dasar yang lebar dengan dinding yang curam,
karena \omterm{def:g8:ice-on-earth:glacier}{gletser} melebarkan dan mendalamkan \omterm{def:g1:landscapes-around-us:landform}{lembah} itu selama mengisinya.

\textbf{10.} $1980 - 1080 = \qty{900}{m}$.

\textbf{11.} $900 \div 18 = 50$ tahun: pada tahun \textbf{2070}.

\textbf{12.} Kecepatannya berubah dari dasawarsa ke dasawarsa (dari 150
sampai \qty{200}{m}), dan mundurnya di masa depan bergantung pada \omterm{def:g4:clouds-rain-snow:precipitation}{hujan}
\omterm{def:g4:clouds-rain-snow:precipitation}{salju} dan suhu musim panas pada tahun-tahun mendatang: laju pada masa
lalu bukanlah sebuah prakiraan.
\end{solution}
